% TOP_PROBLEM: {"problemId": "problem.graph-reconstruction-conjecture", "canonicalTitle": "Graph reconstruction conjecture", "releaseRank": 335, "primaryDomain": "combinatorics_discrete_geometry", "primaryDomainLabel": "Combinatorics & discrete geometry", "displayStatus": "open_with_solved_subcases", "releaseStatus": "open_with_solved_subcases"}
% !TeX program = lualatex
% Definition and sources only; this file is not an LLM solution attempt.
\documentclass[11pt]{article}
\usepackage[margin=25mm]{geometry}
\usepackage{fontspec}
\setmainfont{Latin Modern Roman}
\usepackage{amsmath,amssymb,mathtools}
\usepackage{unicode-math}
\setmathfont{Latin Modern Math}
\usepackage{xurl,hyperref}
\hypersetup{colorlinks=true,urlcolor=blue}
\setcounter{secnumdepth}{0}
\setlength{\parindent}{0pt}
\setlength{\parskip}{5pt}
\setlength{\emergencystretch}{3em}
\begin{document}
\section*{\texorpdfstring{Graph reconstruction conjecture}{Graph reconstruction conjecture}}\label{335}
\subsection{Definitions and mathematical statement}
For a finite simple undirected graph $G=(V,E)$, its \emph{deck} is the multiset of unlabelled isomorphism classes
\[
 \mathcal D(G)=\bigl\{\!\bigl\{[G-v]:v\in V\bigr\}\!\bigr\},
\]
where $G-v$ is the induced graph obtained by deleting $v$ and all incident edges. The reconstruction conjecture asserts that for graphs $G,H$ on at least three vertices,
\[
 \mathcal D(G)=\mathcal D(H)\quad\Longrightarrow\quad G\cong H.
\]
Multiplicity matters: equal-looking vertex-deleted graphs can occur several times in the deck. The vertices removed are not labelled or matched in advance. Edge-deletion reconstruction and reconstruction from only some cards are different questions.
\par\smallskip\textit{Formulation and concept references:} \href{https://www.sciencedirect.com/science/article/abs/pii/S0012365X26000828}{[S1]}.

\subsection{Short English statement}
Can every graph with at least three vertices be recovered up to isomorphism from the multiset of graphs obtained by deleting one vertex?
\subsection{Sources}
\begin{itemize}
\item[S1]Yaxin Qi, Discrete Mathematics 349 (2026), 115058.
\url{https://www.sciencedirect.com/science/article/abs/pii/S0012365X26000828}.
\textit{Formulation source.}
\end{itemize}
\end{document}
